% !TEX root = ../main.tex
\chapter{Conclusions and Future Perspectives}
\label{chap:conclusions}

In this dissertation, a comprehensive Physics-Informed Neural Network (PINN) framework was developed, implemented, and validated for forward field reconstruction and inverse constitutive parameter discovery in complex viscoelastic fluid flows within an enclosed four-roll mill cavity.

\section{Summary of Key Findings and Contributions}
\label{sec:conclusions_summary}

The main scientific and computational findings of this work are summarized below:
\begin{enumerate}
    \item \textbf{Staged Decoupled Optimization Strategy}: 
    Addressing the gradient interference reported in coupled Navier--Stokes and constitutive systems, we demonstrated the effectiveness of a two-stage training schedule: Phase 1 (kinematics and rheology) followed by Phase 2 (hydrodynamic pressure recovery). Retaining an active, trainable stream function $\psi$ during Phase 2 was shown to be necessary to compensate for irrotational velocity components that scalar pressure gradients cannot balance.

    \item \textbf{High-Precision Forward Field Reconstruction}: 
    Validated against dense 125k-node finite element COMSOL solutions for the Oldroyd-B fluid, the PINN achieved sub-percent relative $L_2$ errors on velocity fields ($0.98\%$ for $u$, $0.93\%$ for $v$) and accurately captured localized extra-stress peaks in roller gaps ($1.21\%$ on $\tau_{xy}$, $1.80\%$ on $\tau_{xx}$). Gauge-pressure calibration successfully resolved the hydrostatic offset, yielding a $3.42\%$ relative error on pressure.

    \item \textbf{Inverse Rheological Parameter Discovery}: 
    In inverse discovery benchmarks operating without internal stress measurements, the framework accurately identified physical relaxation times $\lambda$ and polymeric viscosities $\mu_p$ (errors $<1.3\%$ for Oldroyd-B). In the non-linear Giesekus model, the network simultaneously discovered mobility $\alpha$, relaxation time $\lambda$, and viscosity $\mu_p$ within $6.7\%$ error even from deliberately biased blind initializations.

    \item \textbf{Collocation Robustness and Inductive Regularization}: 
    A systematic collocation grid study revealed that parameter estimation accuracy remains consistent across a 25-fold variation in node density (from 5k to 125k nodes), demonstrating that sparse collocation point distributions provide a natural inductive bias against overfitting.

    \item \textbf{Preservation of the Elastic Shear Modulus ($G$)}: 
    Across all constitutive models and parameter sweeps, the network preserved the macroscopic elastic plateau modulus $G = \mu_p / \lambda$ with errors under $1.0\%$, showing that the physics-informed loss prioritizes global elastic stiffness before disentangling individual relaxation timescales.

    \item \textbf{Analysis of the PTT Shear-Flow Degeneracy}: 
    We analyzed a physical limitation of boundary-supervised rheometry for the Phan-Thien--Tanner (PTT) model, where simple shear kinematics near rotating solid walls cause mathematical collinearity between $\lambda$ and $\varepsilon$, yielding an effective response equivalent to a linear Oldroyd-B fluid.
\end{enumerate}

\section{Operational Limitations}
\label{sec:conclusions_limitations}

A critical operational limitation of the current formulation is its reliance on polymeric extra-stress boundary data prescribed along rotating roller surfaces ($\mathcal{L}_{\mathrm{roll\_stress}}$) to anchor absolute stress magnitudes in the absence of internal stress sensors. While physically and mathematically grounded for rotating cylinders, non-intrusively measuring surface polymeric stresses in practical laboratory apparatuses remains experimentally challenging.

\section{Future Perspectives and Research Directions}
\label{sec:conclusions_future_work}

Building upon the findings of this dissertation, several research trajectories warrant future investigation:
\begin{itemize}
    \item \textbf{Integration with Experimental $\mu$-PIV Velocimetry}: 
    Assimilating realistic, noisy Particle Image Velocimetry (PIV) or particle tracking data from physical microfluidic four-roll mill prototypes into the data loss $\mathcal{L}_{\mathrm{data}}$ to evaluate noise tolerance and real-world applicability.
    \item \textbf{Boundary Stress Decoupling via Extended Sensors}: 
    Investigating alternative boundary supervision strategies (such as total wall torque measurements or flow birefringence measurements) to eliminate the requirement for local roller polymeric stress values.
    \item \textbf{Three-Dimensional and Unsteady Viscoelastic Extensions}: 
    Extending the continuous neural surrogate to three-dimensional cavities and time-dependent flows to capture purely elastic 3D flow bifurcations and time-dependent instabilities.
    \item \textbf{Self-Adaptive and Curriculum Optimization}: 
    Incorporating self-adaptive loss weighting schemes (e.g., dynamic weight averaging, gradient-pathology mitigation) and Weissenberg-number continuation curricula to enable stable training at ultra-high elasticities.
\end{itemize}

% ============================================================
% PLACEHOLDER NOTE FOR FUTURE EXPANSION
% ============================================================
% [NOTE: This concluding chapter provides the formal structural summary. 
% Additional domain-specific discussions or experimental comparisons can be inserted here.]
